% Generated draft: requires mathematical and source review.
\documentclass[11pt]{article}
\usepackage[margin=25mm]{geometry}
\usepackage{fontspec}
\setmainfont{Latin Modern Roman}
\usepackage{amsmath,amssymb,mathtools,unicode-math}
\setmathfont{Latin Modern Math}
\usepackage{enumitem,needspace,xurl,hyperref}
\setlength{\parindent}{0pt}
\setlength{\parskip}{5pt}
\setlength{\emergencystretch}{3em}
\newcommand{\N}{\mathbb{N}}
\newcommand{\Z}{\mathbb{Z}}
\newcommand{\Q}{\mathbb{Q}}
\newcommand{\R}{\mathbb{R}}
\newcommand{\C}{\mathbb{C}}
\newcommand{\F}{\mathbb{F}}
\newcommand{\A}{\mathbb{A}}
\newcommand{\PP}{\mathbb{P}}
\newcommand{\E}{\mathbb{E}}
\newcommand{\eps}{\varepsilon}
\newcommand{\dd}{\,\mathrm{d}}
\begin{document}
\section*{UnsolvedMath ID: 30002687}
\textbf{Source problem code:} catalogue entry
\subsection{Definitions and mathematical statement}
Let \((M,g)=(G/K,g)\) be a noncompact, nonflat homogeneous Einstein space, with base point \(o=eK\). In the Riemannian interpretation used here, \(G\) acts transitively by isometries, \(K\) is the stabilizer of \(o\), and
\[
\operatorname{Ric}_g=\lambda g
\]
for a constant \(\lambda\in\R\); nonflat means that the Riemann curvature tensor is not identically zero. No sign condition on \(\lambda\) is supplied.

Assume the stated global Levi decomposition
\[
G=G_1\ltimes G_2,\qquad G_2=AN,
\]
where \(G_1\) is a semisimple Levi factor and \(G_2\) is the solvable radical. The source does not define \(A\) or \(N\), or specify additional properties of their product; the notation \(AN\) is therefore retained without imposing further hypotheses.

Define the solvable-factor orbit and its induced metric by
\[
S=G_2\cdot o\cong G_2/(G_2\cap K),\qquad g_S=\iota_S^*g,
\]
where \(\iota_S:S\longrightarrow M\) is the orbit immersion. This gives a precise induced geometry even when \(G_2\) itself cannot be identified with a subspace of \(M\). Such an identification via \(s\mapsto sK\) requires \(G_2\cap K=\{e\}\), a condition absent from the supplied statement.

For a product decomposition \(G_1=H_1\cdots H_k\) into simple Lie subgroups, let \(O_i=H_i\cdot o\). Write \(\mathfrak h_i\) for the Lie algebra of \(H_i\), and define
\[
X_o^*=\left.\frac{\mathrm d}{\mathrm dt}\right|_{t=0}\exp(tX)\cdot o,
\qquad
V_i=T_oO_i=\{X_o^*:X\in\mathfrak h_i\}.
\]

The supplied problem asks whether the following assertions always hold for the stated data:
\begin{enumerate}
\item The induced metric on the solvable-factor orbit is Einstein: there exists a constant \(\mu\in\R\) such that
\[
\operatorname{Ric}_{g_S}=\mu g_S.
\]
No equality between \(\mu\) and the ambient Einstein constant \(\lambda\) is asserted.
\item Whenever \(G_1=H_1\cdots H_k\) is a product of simple subgroups, their orbits through \(o\) have pairwise orthogonal tangent spaces at \(o\):
\[
\forall\,1\leq i<j\leq k\quad
\forall\,X\in\mathfrak h_i\quad
\forall\,Y\in\mathfrak h_j,\qquad
 g_o(X_o^*,Y_o^*)=0.
\]
\end{enumerate}

These are conditional formulations requiring source clarification. The source does not specify whether additional restrictions on \(A,N,K\), or the choice of Levi factor are intended. The first assertion uses the orbit interpretation of the phrase ``\(G_2\) as a subspace.'' The second uses orthogonality at the common base point; no assertion about orbit distributions elsewhere is supplied. Riemannian signature is the conventional interpretation adopted here, but is not explicitly stated in the record.
\subsection{Short English statement}
For a noncompact, nonflat homogeneous Einstein space with the stated Levi decomposition, ask whether the orbit of its solvable radical, with the induced metric, is Einstein. Also ask whether the simple Levi-factor orbits through the base point have pairwise orthogonal tangent spaces there. The meaning of \(G_2=AN\), the identification of \(G_2\) with an orbit, and any additional intended hypotheses require review.
\subsection{Sources}
\begin{itemize}
\item Supplied local UnsolvedMath registry record, permanent ID \(30002687\), title \emph{Einstein Solvable Factors and Orthogonal Levi Orbits}; source problem code \texttt{catalogue entry}, in the Oberwolfach Reports Open Problems collection. The record gives proposed year 2014 and names no proposer.
\item Supplied source citation: \emph{Mini-Workshop: Einstein Metrics, Ricci Solitons and Ricci Flow under Symmetry Assumptions} (2015), Oberwolfach Reports, pp.\ 2529--2568, DOI \texttt{10.4171/owr/2014/45}. \url{https://doi.org/10.4171/owr/2014/45}. The literature-source list labels this reference ``Oberwolfach Report 45/2014.'' The supplied 2015 citation year and 2014 report designation are retained without reconciliation.
\item Supplied literature reference: Bohm--Lafuente, \emph{Non-compact Einstein Manifolds with Symmetry}. \url{https://doi.org/10.1090/jams/1022}. Inclusion here does not attribute the questions to these authors or establish that this work resolves them.
\item Literature assessment dated exactly \texttt{2026-08-21}: the supplied curation labels the problem open and states that these are auxiliary questions about Levi and solvable factors, including orbit orthogonality, for which later progress on the main Alekseevskii conjecture does not provide a general resolution. This is supplied dated curation, not an independently verified literature assessment; its status and cited claims require independent verification.
\end{itemize}
\end{document}
